% Section 2.5: 本章小结

本章围绕本文研究所依赖的关键理论与技术基础展开：2.1 节梳理了大语言模型的 Transformer 自回归生成机制及代码生成任务在语法严格性、逻辑严密性与语义可执行性等方面的特点与挑战；2.2 节对比全量微调的局限性，重点介绍了本文核心训练侧技术 LoRA 通过低秩矩阵 $W' = W + BA$ 在冻结主干的同时实现轻量化任务适配的原理与优势；2.3 节阐述了链式思维方法的两条主线，即不依赖示例的 Zero-shot CoT 与基于示例引导的 Few-shot CoT，及其在代码生成场景中将「推理」与「编码」显式分离的工程价值；2.4 节明确了以可执行单测为基础、Pass@k 无偏估计 $\text{pass@k} = 1 - \binom{n-c}{k} / \binom{n}{k}$ 为度量的评测协议。上述四节共同构成第三至第五章「数据—方法—实验」的概念底座：训练侧 LoRA 决定参数侧的可塑性边界，推理侧 CoT 决定生成阶段的结构化约束，Pass@k 给出对二者协同效果的可比量化口径，三者在第四章的析因实验矩阵中以正交因子的形式接受系统性检验。